

header-includes:
  - \usepackage{helvet}
  - \renewcommand*\familydefault{\sfdefault}
  - \usepackage{setspace}\doublespacing
  - \usepackage[left]{lineno}
  - \linenumbers
  - \usepackage{dcolumn}
  - \usepackage{caption}
  - \usepackage{float}
  - \usepackage{afterpage}
  - \usepackage{siunitx}
  - \usepackage{amsmath}
% fontfamily: mathpazo
% fontsize: 11pt
% geometry: margin = 1.0in
keywords: Malaria, Transmission dynamics, simulation, models
# bibliography: ./library.bib
# csl: ./nature-conservation.csl

% abstract: |
%   Recently documented insect declines suggest that agricultural expansion and productivity-focused intensification are a putative driver for the observed declines. However, this hypothesis has yet to be formally tested due to a lack of historical data on agricultural land-use change and insect taxa. Using large, spatiotemporal sets of bumble bee and agricultural records, we show that agricultural intensification, specifically an increase cropland extent and insecticides, as well as a decrease in crop richness and perennial habitat, are associated with a decrease in the occurrence of bumble bee species in the agriculturally intensive Midwest. Yet, we found positive associations between bumble bees and crop richness and pasture, even in areas where agriculture is spatially extensive. Our findings suggest that insect conservation and agricultural production may be compatible, with increasing in on-farm and landscape-level crop diversity, pasture, and reduction in insecticides predicted to have positive effects on bumble bees.  